\section{后训练革命：RLVR 让“可验证奖励”成为主路径}

\subsection{RLHF、RLVR 与完整 post-training recipe}

访谈里最清晰的部分之一是后训练配方：mid-training 先给任务结构，RLVR 用可验证目标做强化学习，RLHF 最后做风格和可用性收敛。这个顺序解释了为什么很多新模型在数学与代码上进步快，同时在“说话方式”上保持可控。

\begin{importantbox}{RLVR 为什么在 2025--2026 爆发}
\begin{itemize}[leftmargin=1.5em]
    \item 数学、代码等任务可自动验对错，奖励信号清晰；
    \item 奖励函数不依赖“模拟偏好”的 reward model，降低 reward hacking 风险；
    \item 工程上可并行 rollout，便于持续扩展训练吞吐。
\end{itemize}
\end{importantbox}

\begin{figure}[H]
\centering
\includegraphics[width=0.88\textwidth]{frame-02.jpg}
\caption{访谈中对 RLVR 与后训练配方的集中讨论\protect\footnotemark}
\end{figure}
\footnotetext{视频画面时间区间：01:47:10--01:47:20。}

\subsection{Actor--Learner 架构与基础设施代价}

Nathan 把 RL 训练解释为两类计算节点协同：actor 负责生成与采样，learner 负责梯度更新。前者可地理分散，后者需要紧密互联。这个结构与传统 pre-training 的单一大集群模式不同，因此调度、通信、容错、数据回流都更复杂。

\begin{knowledgebox}{为什么 RL 工程经常“看起来比预训练更乱”}
\begin{itemize}[leftmargin=1.5em]
    \item 数据是在线生成的，不是静态语料直接遍历；
    \item 回报信号可能延迟出现，调参反馈更慢；
    \item 失败案例分析既要看策略轨迹，也要看系统日志；
    \item 推理与训练耦合，任何一侧抖动都可能拖垮整体效率。
\end{itemize}
\end{knowledgebox}

\begin{warningbox}{RLVR 不是“自动生成新知识”机器}
访谈也明确保留了怀疑：RL 更像把模型已有能力挖出来并重排，而不是凭空注入新知识。若预训练语料对某类问题覆盖不足，RLVR 可以提升过程质量，但很难凭空制造缺失事实。这个边界对路线规划很关键。
\end{warningbox}

\subsection{本章小结}

后训练阶段已从“微调附属环节”升级为主战场。RLVR 的爆发来自任务可验证性与工程可扩展性的结合，但它并不替代高质量预训练数据。

